\documentclass[UTF8,11pt,oneside]{ctexart}
\usepackage[margin=2.2cm]{geometry}
\usepackage{amsmath,amssymb,bm,mathtools}
\usepackage{booktabs,array,tabularx}
\usepackage{enumitem}
\usepackage{tikz}
\usetikzlibrary{arrows.meta,positioning,fit,calc}
\usepackage{placeins}
\usepackage{hyperref}
\hypersetup{colorlinks=true,linkcolor=blue,urlcolor=blue,citecolor=blue}
\setlength{\parskip}{0.15em}
\setlength{\emergencystretch}{3em}
\setlength{\parindent}{2em}
\setlist[enumerate]{itemsep=.2em,topsep=.2em}
\setlist[itemize]{itemsep=0pt,topsep=0pt}
\newcommand{\R}{\mathbb{R}}
\newcommand{\relu}{\operatorname{ReLU}}
\newcommand{\softmax}{\operatorname{Softmax}}
\newcommand{\Lcal}{\mathcal{L}}
\title{作业 1：反向传播与神经网络}
\author{深度学习（84100343-0）\\清华大学}
\date{2026 年秋季}
\begin{document}
\pagestyle{plain}
\begin{center}
{\LARGE\bfseries 作业 1：反向传播与神经网络}\par\medskip
深度学习（84100343-0）\quad 清华大学\par
2026 年秋季\quad {\small 中文译本}
\end{center}

\section*{引言}
\textbf{构建一个由专家组成的小型语言模型。}莎士比亚文本和 Python 代码遵循不同的模式，但都可以通过预测下一个符号来建模。是否每个 token 都必须使用同一个前馈网络？你将构建一个混合专家（MoE）模型，并研究它如何在控制每个 token 计算量的同时提升模型容量。

MoE 层存储若干个 MLP，称为\emph{专家}。一个经过学习的\emph{路由器}会为每个 token 选择少数几个专家，并组合它们的输出，从而将总容量与实际参与的计算分开。专家通过联合训练学习各自的作用，而非人工指定。Switch Transformers 在大规模场景探索了这一原理 [1]。Mixtral 8x7B 在语言、代码和多语言任务中为每个 token 从八个专家中选择两个 [2]。DeepSeek-V3 报告总计 6710 亿参数、每个 token 激活 370 亿参数 [3]。本作业将在课程作业规模下研究这一机制。

\begin{figure}[ht]
\centering
\begin{tikzpicture}[>=Latex, node distance=7mm, every node/.style={font=\small}, box/.style={draw,rounded corners,minimum height=7mm,inner sep=4pt}]
\node[box] (x) {Token 状态 $\bm{x}$};
\node[box,fill=gray!20,right=of x] (r) {路由器 / Top-2};
\node[box,fill=blue!15,right=13mm of r] (e1) {专家 1};
\node[box,right=of e1] (e2) {专家 2};
\node[box,fill=blue!15,right=of e2] (e3) {专家 3};
\node[box,right=of e3] (e4) {专家 4};
\node[box,fill=orange!30,right=12mm of e4] (o) {加权输出};
\draw[->] (x)--(r);
\draw[->] (r.east) -- ++(4mm,0) |- (e1.north);
\draw[->] (r.east) -- ++(4mm,0) |- (e3.north);
\draw[->] (e1.south) -- ++(0,-5mm) -| (o.south west);
\draw[->] (e3.south) -- ++(0,-5mm) -| (o.south west);
\node[font=\scriptsize,below=1mm of e2] {不执行};
\node[font=\scriptsize,below=1mm of e4] {不执行};
\node[font=\scriptsize,below=7mm of e3] {存储多个专家；每个 token 仅激活少数专家};
\end{tikzpicture}
\caption{一个 token 激活所存储专家中的一个小子集。参照 Mixtral [2] 图 1 的路由模式原创绘制的示意图；此处为便于说明仅画出四个专家。}
\end{figure}

我们的下一字节预测器包含三个步骤：推导基本模块、用 NumPy 验证它们，以及测试一个小型因果语言模型。专家采用常见的全连接（FC）层和 ReLU；因果上下文模块已提供。请遵守以下要求：
\begin{itemize}
\item 在书面报告（PDF）中逐项作答，语言可为英文或中文，并提供清晰且带注释的代码。不得分享或提交他人的作业。
\item \textbf{[手写]}题要求独立手写推导。请将作业拍照或扫描后插入 PDF。不得使用 AI 生成推导；可用计算器进行数值近似。
\item 未标记的题目要求自行编写核心代码；AI 可以解释一般语法或环境错误，但不能提供该核心代码。[AI 探索]题鼓励借助 AI 探索。请披露所获协助及讨论内容；不得捏造结果。
\end{itemize}
总分：100 分（手写 35 分，核心编程 50 分，探索 15 分）。代码和数据随本 PDF 提供；请在发布包目录下运行命令。

\begin{center}\small 深度学习（2026 年秋季），清华大学。\end{center}

\clearpage
\section{第一部分：反向传播（35 分）}
\subsection{引言}
稀疏 MoE 层为每个 token 表示选择少数几个前馈专家 [1,2]。其组成部分都很常见：全连接层、ReLU、Softmax、加权求和和残差连接。我们先研究一个小型可微示例。局部卷积提供上下文表示；第三部分会用已提供的因果上下文模块替换此编码器。

\subsection{网络细节}
给定输入 $\bm{s}\in\R^4$，共享卷积核 $\bm{q}\in\R^3$ 产生
\begin{equation}
a_j=\sum_{t=1}^{3}q_t s_{j+t-1},\qquad h_j=\relu(a_j),\qquad j=1,2.\tag{1}
\end{equation}
不使用填充，步幅为 1，运算为互相关。该编码器给出 $D=2$；下游模块按一般 $D$ 编写。取 $\relu'(0)=0$。对于 $E$ 个专家、隐藏宽度 $H$ 和 $K$ 个预测类别，定义
\begin{align}
\bm{v}_e&=W_{1,e}\bm{h}+\bm{b}_{1,e}, & \bm{u}_e&=\relu(\bm{v}_e), & \bm{f}_e&=W_{2,e}\bm{u}_e+\bm{b}_{2,e}, \tag{2}\\
\bm{g}&=\softmax(G\bm{h}), & \bm{r}&=\bm{h}+\sum_{e=1}^{E}g_e\bm{f}_e, & \hat{\bm{y}}&=\softmax(C\bm{r}+\bm{b}_c). \tag{3}
\end{align}
其中 $W_{1,e}\in\R^{H\times D}$、$W_{2,e}\in\R^{D\times H}$、$G\in\R^{E\times D}$、$C\in\R^{K\times D}$。所有偏置的维度与其对应输出相同；路由器没有偏置。对大小为 $m$ 的小批量，使用
\begin{equation}
\Lcal=-\frac{1}{m}\sum_{i=1}^{m}\sum_{k=1}^{K}y_k^{(i)}\log\hat{y}_k^{(i)}.\tag{4}
\end{equation}

\begin{figure}[htbp]
\centering
\resizebox{\linewidth}{!}{%
\begin{tikzpicture}[>=Latex,every node/.style={font=\small},box/.style={draw,rounded corners,align=center,minimum height=6mm,inner sep=3pt}]
\node[box] (s) at (0,0) {输入 $\bm{s}$};
\node[box,fill=green!15] (conv) at (1.2,0) {卷积};
\node[box,fill=red!15] (h) at (2.4,0) {ReLU};
\node[box,fill=blue!15] (w11) at (4.2,.9) {FC $W_{1,1}$};
\node[box,fill=red!15] (a1) at (5.8,.9) {ReLU};
\node[box,fill=blue!15] (w21) at (7.4,.9) {FC $W_{2,1}$};
\node[box,fill=blue!15] (w12) at (4.2,-.9) {FC $W_{1,2}$};
\node[box,fill=red!15] (a2) at (5.8,-.9) {ReLU};
\node[box,fill=blue!15] (w22) at (7.4,-.9) {FC $W_{2,2}$};
\node[box,fill=blue!15] (g) at (4.2,-2) {FC $G$};
\node[box,fill=gray!20] (gate) at (6.2,-2) {Softmax};
\node[box,fill=orange!25] (sum) at (9.5,0) {加权求和};
\node[draw,circle,inner sep=2pt] (plus) at (10.9,0) {$+$};
\node[box,fill=blue!15] (fc) at (12,0) {FC $C$};
\node[box,fill=gray!20] (sm) at (13.6,0) {Softmax};
\node[box] (pred) at (15.1,0) {预测};
\draw[->] (s)--(conv);\draw[->] (conv)--(h);
\draw (h.east)--node[above]{$\bm{h}$}(3.2,0);
\draw[->] (3.2,0)|-(w11.west);\draw[->] (3.2,0)|-(w12.west);\draw[->] (3.2,0)|-(g.west);
\draw[->] (w11)--(a1);\draw[->] (a1)--(w21);
\draw[->] (w12)--(a2);\draw[->] (a2)--(w22);
\draw[->] (w21.east)--node[above]{$\bm{f}_1$}(sum.north west);
\draw[->] (w22.east)--node[below]{$\bm{f}_2$}(sum.south west);
\draw[->] (g)--(gate);
\draw[->] (gate.east)--node[below]{$g_1,g_2$}(9.5,-2)--(sum.south);
\draw[->] (3.2,0)--(3.2,1.8)--node[above]{残差连接 $\bm{h}$}(10.9,1.8)--(plus.north);
\draw[->] (sum)--(plus);\draw[->] (plus)--node[above]{$\bm{r}$}(fc);\draw[->] (fc)--(sm);\draw[->] (sm)--(pred);
\node[above=1mm] at (5.8,1.2) {专家 1};\node[below=1mm] at (5.8,-1.2) {专家 2};
\end{tikzpicture}%
}
\caption{手写推导示例网络。蓝、红、灰、橙色框分别表示全连接层、ReLU、Softmax 和加权求和；圆形表示加法。两个专家均被选中。}
\end{figure}

数值计算使用 $m=1,D=H=E=2,K=3$，参数如下。未列出的专家偏置均为零：
\begin{align*}
\bm{s}&=(2,1,-1,3)^\top,&\bm{q}&=(1,\tfrac12,\tfrac12)^\top,&\bm{y}&=(0,1,0)^\top,\\
W_{1,1}&=\begin{pmatrix}1&0\\0&-1\end{pmatrix},&W_{1,2}&=\begin{pmatrix}0&1\\-1&0\end{pmatrix},&\bm{b}_{1,1}&=\begin{pmatrix}0\\1/2\end{pmatrix},&\bm{b}_{1,2}&=\begin{pmatrix}0\\1\end{pmatrix},\\
W_{2,1}&=\begin{pmatrix}1/5&0\\0&1/5\end{pmatrix},&W_{2,2}&=\begin{pmatrix}0&1/5\\1/5&0\end{pmatrix},&\bm{b}_{2,1}&=\begin{pmatrix}0\\1/2\end{pmatrix},&\bm{b}_{2,2}&=\begin{pmatrix}1/2\\0\end{pmatrix},\\
G&=\begin{pmatrix}1/4&-1/4\\0&0\end{pmatrix},&C&=\begin{pmatrix}1/5&0\\0&1/5\\1/5&1/5\end{pmatrix},&\bm{b}_c&=(-0.49,-0.49,-0.98)^\top.\tag{5}
\end{align*}

\subsection{题目}
请展示必要的推导和变量维度。强烈建议使用向量形式；中间伴随量可以重复使用。
\begin{enumerate}
\item \textbf{[手写]} 计算式 (5) 对应的 $\bm{h}$、两个专家的输出、$\bm{g}$、$\bm{r}$、预测值和损失。（6 分）
\item \textbf{[手写]} 推导 Softmax 雅可比矩阵，以及交叉熵对预测 logits 的梯度。说明小批量平均如何进入梯度。（6 分）
\item \textbf{[手写]} 对一般的 $m,D,H,E,K$，推导 $W_{1,e},\bm{b}_{1,e},W_{2,e},\bm{b}_{2,e},G,C,\bm{b}_c$ 以及每个输入 $\bm{h}$ 的梯度。明确包含专家、路由器和残差路径。仅数值计算 $\partial\Lcal/\partial G_{1,1}$。（16 分）
\item \textbf{[手写]} 通过共享卷积反向传播。可使用经核验的上游梯度 $\nabla_{\bm{h}}\Lcal=(349,-173)^\top/2400$。计算 $\nabla_{\bm{q}}\Lcal$、$\partial\Lcal/\partial s_2$，并给出一次学习率为 0.1 的卷积核更新。为所有可训练参数给出小批量 SGD 伪代码。（7 分）
\end{enumerate}

\section{第二部分：实现专家（20 分）}
\subsection{引言}
使用 NumPy 实现式 (2)--(4) 中的模块。不允许使用自动微分。完成 \texttt{NumPy/model.py}；数据、初始化、中心差分检查器和训练循环均已提供。批量样本按行存储；参数数组遵循上文的列向量约定。
\subsection{稀疏路由与梯度验证}
对于选出的 $k$ 个专家，令 $S(\bm{h})$ 表示路由 logits 中最大的 $k$ 个所对应的索引，并将式 (3) 中的门控替换为
\begin{equation}
g_e=\begin{cases}\displaystyle\frac{\exp((G\bm{h})_e)}{\sum_{j\in S(\bm{h})}\exp((G\bm{h})_j)},&e\in S(\bm{h}),\\[6pt]0,&e\notin S(\bm{h}).\end{cases}\tag{6}
\end{equation}
选中集合在局部固定。检查器会提供远离 Top-k 并列情况和 ReLU 折点的掩码；无需对离散选择求导。检查器使用 $\epsilon=10^{-6}$，并报告最大归一化误差
\begin{equation}
\max_j\frac{|a_j-n_j|}{\max(1,|a_j|+|n_j|)},\qquad n_j=\frac{\Lcal(\theta+\epsilon e_j)-\Lcal(\theta-\epsilon e_j)}{2\epsilon}.\tag{7}
\end{equation}
\subsection{题目}
\begin{enumerate}
\item 实现前向与反向传播，包括预测头和共享输入梯度。复现第一部分，并支持所提供的三专家 Top-2 情形。（10 分）
\item 运行 \texttt{python}\allowbreak\ \texttt{NumPy/run.py}\allowbreak\ \texttt{--check}。报告每类参数和输入的误差，目标是误差低于 $10^{-5}$。诊断所提供的错误输入梯度变体，并指出缺少了哪条计算路径。（6 分）
\item 在所提供的同一语料库的 128 个样本批次上运行 \texttt{python}\allowbreak\ \texttt{NumPy/run.py}\allowbreak\ \texttt{--train}。使用 400 次 SGD 更新和 0.2 的学习率；绘制损失曲线并展示三个预测。上下文特征固定。不要求进行超参数搜索或达到泛化目标。（4 分）
\end{enumerate}

\section{第三部分：小型 MoE 语言模型（45 分）}
\subsection{数据集}
在 Tiny Shakespeare [4] 与 CPython 源代码 [5] 的固定混合数据上预测下一个字节。词表由全部 256 种可能字节组成。软件包提供了分词后的数据和固定划分：Shakespeare 使用连续的 80/10/10 百分比划分；代码按文件划分。窗口不会跨越文件或数据划分边界。训练时两个数据域等权采样。域标签仅用于评估，绝不作为路由器输入。
\subsection{网络细节}
所提供的双块因果模型维度为 $D=64$，有四个注意力头，上下文长度为 $T=64$。因果上下文模块只使用当前及之前的字节。嵌入、归一化、注意力、残差包装器、优化器和负载均衡均已提供；不评估其内部实现。只需完成 \texttt{MoE/model.py} 中标记的核心方法。本部分允许使用 PyTorch 自动微分。

每个专家使用 ReLU，将维度 $64\to128\to64$。一个块的结构是 $\bm{x}+\operatorname{MoE}(\operatorname{LN}(\bm{x}))$，并采用式 (6)，其中 $E=16,k=2$。第一部分为手算省略了归一化包装器。这是用于教学的小规模 token 级 FFN 路由模型，并非对预训练大型模型的复现。

\begin{figure}[htbp]
\centering
\resizebox{\linewidth}{!}{%
\begin{tikzpicture}[>=Latex,every node/.style={font=\small},box/.style={draw,rounded corners,align=center,minimum height=7mm,inner sep=4pt}]
\node[box] (bytes) at (0,0) {字节};
\node[box,fill=green!15] (emb) at (1.4,0) {嵌入};
\node[box,fill=green!15] (ctx) at (2.9,0) {上下文};
\node[box,fill=green!15] (ln1) at (4.5,0) {归一化};
\node[box,fill=blue!15] (moe) at (6.2,0) {MoE FFN};
\node[draw,circle,inner sep=2pt] (add) at (7.7,0) {$+$};
\node[box,fill=green!15] (norm) at (9.1,0) {归一化};
\node[box,fill=blue!15] (head) at (10.8,0) {FC 输出头};
\node[box,fill=gray!20] (sm) at (12.6,0) {Softmax};
\node[box] (next) at (14.5,0) {下一字节};
\foreach \a/\b in {bytes/emb,emb/ctx,ctx/ln1,ln1/moe,moe/add,add/norm,norm/head,head/sm,sm/next}{\draw[->](\a)--(\b);}
\draw[->] (3.7,0)--(3.7,1)--node[above]{$\bm{x}$}(7.7,1)--(add.north);
\draw (2,-.55)--(2,-.8)--(8.1,-.8)--(8.1,-.55);
\node[below] at (5.05,-.8) {两个块，参数各自独立};
\begin{scope}[yshift=-8mm]
\node[box] (ln) at (0,-2.9) {LN($\bm{x}$)};
\node[box,fill=blue!15] (g) at (1.8,-2.9) {FC $G$};
\node[box,fill=gray!20] (top) at (3.4,-2.9) {Top-2};
\node[box,fill=blue!15,text width=37mm] (ea) at (7,-2.1) {专家 $a$：FC / ReLU / FC};
\node[box,fill=blue!15,text width=37mm] (eb) at (7,-3.7) {专家 $b$：FC / ReLU / FC};
\node[box,fill=orange!25] (weighted) at (10.5,-2.9) {加权求和};
\node[box] (out) at (12.5,-2.9) {MoE 输出};
\draw[->] (ln)--(g);\draw[->] (g)--(top);
\draw (ln.north)--(0,-1.6)--node[above]{token 表示}(4.6,-1.6);
\draw[->] (4.6,-1.6)|-(ea.west);\draw[->] (4.6,-1.6)|-(eb.west);
\draw[->] (ea.east)--(weighted.north west);\draw[->] (eb.east)--(weighted.south west);
\draw[->] (top.south)--(3.4,-4.5)--node[below]{选中的索引及归一化门控权重}(10.5,-4.5)--(weighted.south);
\draw[->] (weighted)--(out);
\node[above] at (7,-1.55) {16 个专家；仅执行 2 个};
\node[align=left] at (13.3,-4) {$D=64,\ T=64$\\专家宽度 $=128$};
\end{scope}
\end{tikzpicture}%
}
\caption{训练模型。两个块结构相同，但参数独立。仅被选中的 token--专家对执行专家 MLP。绿色为已提供模块，蓝色为全连接层，灰色为门控或概率。}
\end{figure}

\begin{table}[ht]
\centering
\caption{三种配置；其余模型组件均相同。参数量匹配为近似值（误差在 1\% 以内）。}
\begin{tabularx}{\linewidth}{>{\raggedright\arraybackslash}Xccc}
\toprule
配置 & 专家数 & 选中数 & FFN 隐藏宽度\\
\midrule
计算量匹配的稠密模型 & 1 & 1 & 264\\
稀疏 MoE & 16 & 2 & 每个专家 128\\
参数量匹配的稠密模型 & 1 & 1 & 2056\\
\bottomrule
\end{tabularx}
\end{table}

\FloatBarrier
\subsection{训练与评估}
使用随机种子 416、3000 次更新，以及由 32 个长度为 64 的序列组成的批次。保留提供的优化器、学习率计划、数据顺序和检查点规则：选取验证交叉熵最低的检查点。辅助负载均衡项的系数为 0.01，无需手工推导。评估时使用留出的、互不重叠且长度完整的数据块。分别报告两个数据域，并以
\begin{equation}
\Lcal_{\mathrm{test}}=\tfrac12(\Lcal_{\mathrm{text}}+\Lcal_{\mathrm{code}}),\qquad \mathrm{PPL}=\exp(\Lcal_{\mathrm{test}})\tag{8}
\end{equation}
作为主要指标。只有在词元化方式和评估协议相同的前提下，字节级困惑度才可比较。本比较固定更新次数和数据顺序，不固定墙钟时间或收敛程度。评分取决于正确性和分析，而非 MoE 是否胜出或是否达到绝对速度/PPL 目标；不同设备间可能存在小的数值差异。

\subsection{题目}
\begin{enumerate}
\item 实现专家前向计算、路由器和所选 token 的分发。先运行提供的 CPU 检查 \texttt{python}\allowbreak\ \texttt{MoE/diagnostics.py}\allowbreak\ \texttt{--check}；然后执行 \texttt{python}\allowbreak\ \texttt{MoE/train.py}\allowbreak\ \texttt{--mode all}。报告全部三种配置的训练/验证曲线和测试指标。不要先计算所有专家再屏蔽其输出。（10 分）
\item 在同一设备上比较质量、总参数量、每 token 的矩阵乘加（MAC）次数，以及完整批次前向延迟（$B=32,T=64$）。MAC 统计要包括注意力、路由器、专家和输出头；计时要包括分发开销。运行 \texttt{python MoE/diagnostics.py}：在最后一层同一位置上，不同序列之间交换成对路由，同时保持专家--token 数量不变。该程序会检查前缀不变性，并报告三个固定打乱种子下的诊断交叉熵 CE。这种依赖批次的扰动不代表语言模型独立运行时的性能；不要据此推断因果训练收益。用一张表和一幅图讨论容量、计算、路由和专家使用情况。仅凭使用量不能证明语义专门化；较低 MAC 也不一定意味着更低延迟。（20 分）
\item \textbf{[AI 探索]} 设计一个小型自主研究。选择一个关于稀疏计算、路由或模型行为且可检验的问题。AI 工具可协助构思、编程和制图。（15 分）
\end{enumerate}
可选的起始方向（也欢迎其他问题）：(a) 在一个 Top-2 检查点上比较推理时 $k=1,2,4$：质量与 MAC 的关系；(b) 固定最后一层选中的专家，比较学习得到的门控权重与相等门控权重（各为 $1/2$）：仅仅选对专家是否足够？

明确提出问题和有意义的对照；围绕必做实验结果，通过少量实验或分析收集证据。解释你的理解及一个局限；不要求结果有所改进。研究报告最多 600 个英文单词或 1000 个中文字，图表最多两幅；请提供可运行代码和 AI 使用声明。评分依据问题（3 分）、设计（4 分）、证据（4 分）、解释（3 分）和披露（1 分）。

\section{提交格式}
提交名为 \texttt{\{StudentID\}\_\{Name\}\_hw1.zip} 的压缩包，其中包含：
\begin{enumerate}
\item 一份逐题作答的 PDF，包括第一部分的手写照片、梯度检查结果、实验比较和探索研究。
\item 第二、三部分可运行的核心代码、配置和精简结果日志。请包含必要的参考文献和协助披露；无需包含完整模型检查点和所提供的数据集。
\end{enumerate}

\section*{AI 使用声明}
参照 ICLR 的做法，请在报告结尾用 3--5 句话披露所用工具/模型、受影响的题目/任务及核验方式，或明确声明未使用 AI。确认遵守以上规则，并对提交内容负责。本声明不计入 600 词/1000 字的限制。

\section*{参考文献}
\footnotesize
\begin{enumerate}[label={[\arabic*]}]
\item W. Fedus, B. Zoph, and N. Shazeer. Switch Transformers: Scaling to Trillion Parameter Models with Simple and Efficient Sparsity. JMLR, 23(120), 2022.
\item A. Q. Jiang et al. Mixtral of Experts. arXiv:2401.04088, 2024.
\item DeepSeek-AI et al. DeepSeek-V3 Technical Report. arXiv:2412.19437, 2025.
\item A. Karpathy. Tiny Shakespeare corpus, char-rnn. \url{https://github.com/karpathy/char-rnn}.
\item Python Software Foundation. CPython 3.10.16 source distribution. \url{https://www.python.org/downloads/release/python-31016/}.
\end{enumerate}
\end{document}
